\section{Complete native multiloop reductions}
\label{sec:complete-multiloop}
The preceding supplied-equation tests begin with an existing IBP pool.
To include the work needed to form that pool, we now give Kira physical
denominators and indexed target integrals and repeat equation generation
in every measured calculation. We also train a separate row model from
the algebraic work caused by choices inside Kira. This connects the
training objective to the intended reduction cost and permits a direct
comparison with Kira's own source-row choices.

\subsection{Physical families and independent targets}
The integral notation remains $I_{\boldsymbol a}$, with
$D_i=q_i^2-m_i^2+\mathrm{i}0$ as defined earlier. We use the following
ordered momentum lists:
\begin{align}
 \boldsymbol q_{\rm sunrise}
 &= (\ell_1,\ell_2,p-\ell_1-\ell_2,p-\ell_1,p-\ell_2),\\
 \boldsymbol q_{\rm box}
 &= (\ell_1,\ell_1-p_1,\ell_1-p_1-p_2,
     \ell_2,\ell_2-p_1-p_2,\ell_2-p_1-p_2-p_3,\nonumber\\
 &\hspace{2em}\ell_1-\ell_2,\ell_1-p_3,\ell_2-p_1),\\
 \boldsymbol q_{\rm vacuum}
 &= (\ell_1,\ell_2,\ell_3,\ell_1-\ell_2,
     \ell_2-\ell_3,\ell_1-\ell_3),\\
 \boldsymbol q_{\rm ladder}
 &= (\ell_1,\ell_1-p,\ell_2,\ell_2-p,\ell_3,\ell_3-p,
     \ell_1-\ell_2,\ell_2-\ell_3,\ell_1-\ell_3).
\end{align}
The first three sunrise entries, seven box entries and eight ladder
entries are physical propagators. Their remaining entries complete the
scalar-product basis and can appear with nonpositive indices. All six
vacuum entries are physical propagators. The first two lists contain
two loop momenta, and the last two contain three. In particular, the
three-loop vacuum and ladder systems couple the loop momenta through
shared propagators.

For each list, we expand the squared momenta in
$\ell_i\cdot\ell_j$ and $\ell_i\cdot p_j$ and check that the resulting
coefficient matrix has rank
\[
 N_{\rm sp}=\frac{L(L+1)}{2}+LE.
\]
Here $E$ is the number of independent external momenta, and $N_{\rm sp}$
is the number of independent scalar products involving loop momenta.
All fifteen family definitions pass this check, and their nonempty
external Gram matrices have nonzero determinants. This prevents a
dependent denominator list from entering the comparison.

All squared masses and external scalar products have fixed rational
values in a common squared-momentum unit, while the dimension $d$
remains symbolic. The reduction coefficients therefore belong to
$\mathbb Q(d)$. For the box, the external momenta satisfy
\[
 p_i^2=0,\qquad
 2p_1\cdot p_2=s,\qquad
 2p_1\cdot p_3=t,\qquad
 2p_2\cdot p_3=-s-t.
\]
The six training cases consist of two sunrise mass profiles
$(1,1,1)$ and $(1,2,3)$ with $p^2=-1,-2$ respectively, two massless
boxes with $(s,t)=(-1,-2),(-2,-3)$, and two vacuum profiles
$(1,1,1,1,1,1)$ and $(1,1,1,2,2,2)$. These profiles specify squared
masses in the displayed propagator order. Validation uses the sunrise
profile $(1,1,2)$ at $p^2=-3$, a box at $(s,t)=(-3,-4)$, and the vacuum
profile $(1,1,1,1,1,2)$.

The six reserved cases use boxes at $(s,t)=(-5,-7),(-7,-11)$,
vacuum profiles $(1,1,2,1,2,2)$ and $(1,2,2,1,2,2)$, and two ladder
instances. One ladder is massless with $p^2=-1$; the other has
$m_1^2=1$, all other squared masses zero, and $p^2=-2$. The ladder
topology appears only in the reserved set. The other reserved cases
test different kinematic or mass conditions within training topologies.
These tests concern rational kinematic points with symbolic $d$ and
do not establish performance for general symbolic multiscale
coefficients.

Targets are individual indexed integrals. We begin with the full-sector
corner, then increase each positive index by one and insert each
available auxiliary numerator once. We repeat this enumeration after
removing one, two and three physical denominators, retain at least two
physical denominators, and remove duplicate index vectors. The reserved
batches contain the first 32 or 128 targets in this fixed order. This
construction selects targets without using learned scores or reverse
construction paths.

\subsection{Seed regions, learned subsets and complete fallback}
For a batch $\mathcal T$, define
\[
 r_{\mathcal T}=\max_{\boldsymbol a\in\mathcal T}
                  \sum_i\max(a_i,0),\qquad
 s_{\mathcal T}=\max_{\boldsymbol a\in\mathcal T}
                  \sum_i\max(-a_i,0),\qquad
 b_{\mathcal T}=\max_{\boldsymbol a\in\mathcal T}
                  \sum_i\max(a_i-1,0).
\]
The last quantity is the largest total number of additional propagator
powers, also called dots. We specify Kira seed bounds
$(r_{\mathcal T}+\delta,s_{\mathcal T}+\delta,b_{\mathcal T}+\delta)$,
where $\delta=1$ defines the conservative reference region and
$\delta=0$ defines the smaller region. Kira still applies its native
symmetry, equation-generation and equation-selection procedures within
these bounds~\cite{Kira3}. The bounds are explicit input choices;
their difference is not a learned improvement to Kira.

The saved seed model scores the finite candidate region using the
ordered sum of the batch targets, with coefficient one for each target.
This sum supplies model context, while Kira still reduces every target
separately. Every sector corner, target and prescribed master that
belongs to the candidate region remains available. An additional
deterministic rule can retain all seeds with at most one dot or one
numerator power in total. The learned policy can add seeds above its
probability threshold, and a separate control applies the same
deterministic rule without model inference.

Kira first checks whether the selected system requires masters outside
the prescribed reference basis. A rejected subset triggers the
documented fallback sequence, and its time remains part of the
measurement. A final successful calculation must produce every
requested target and match its reference coefficient in $\mathbb Q(d)$.
Other failed attempts also move through the prescribed fallback sequence. A coefficient mismatch makes the complete workflow unsuccessful. Splitting a batch into groups of
32 targets bounds each calculation, but it also repeats equation
generation and process startup; the timing comparison includes both
effects.

\subsection{Training row choices on their algebraic consequences}
The new row model chooses a source among equations with the same
leading integral. It preserves Kira's native pivot and group order, and
Kira carries out every resulting row operation exactly. At state $S$,
the allowed action $q$ selects one such source row, which removes the
leading integral from the other rows in that group. The model therefore
acts on the changing native matrix through a restricted, explicitly
defined action set.

The first training objective measures the immediate symbolic insertion
time for each source. We clone the destination rows, apply the exact
native update, and average two rotated measurement orders while
excluding clone time. This measures local arithmetic cost, but it does
not include the equations created by later eliminations. We therefore
also measure a continuation cost,
\begin{equation}
 C(S,q)=\sum_{j\in\mathcal U(S,q)}(1+2n_j).
 \label{eq:continuation-cost}
\end{equation}
Here $\mathcal U(S,q)$ is the sequence of committed forward insertions
after choosing $q$ and then following Kira's native choices, and $n_j$
is the number of nonzero coefficients in the source of insertion $j$.
The count assigns one division and two coefficient operations per
source entry. It is an algebraic operation estimate tied to measured
native insertions; it does not count Fermat's internal polynomial
operations or back substitution.

We choose up to four evenly spaced eligible states from the saved
local-cost sample in each development case, which contains at most
100 measured groups. We prefer groups with at least three source
choices and use smaller groups when no such group is available. Each alternative receives a separate complete Kira continuation.
The recorded pivot and all candidate features must match the original
state, and the final master coefficients must match the reference
reduction before its cost can enter training. The resulting labels
measure a choice followed by native continuation. They do not prove a
globally minimum reduction path.

Each candidate has eighteen features describing row length, coefficient
text length and operations, overlap with destination columns, estimated
new entries, and its position relative to Kira's preferred source.
A linear fit to log continuation cost retained the native source in all
nine validation states. We therefore also fit a small regression tree~\cite{CART1984},
which can represent conditions involving several features.

For this tree, the target is the normalized change in measured
continuation cost,
\[
 y_q=\frac{C(S,q)-C(S,0)}
           {\max\{1,\max_{q'}C(S,q')-\min_{q'}C(S,q')\}},
 \qquad
 \widehat y_q=h(\boldsymbol x_q-\boldsymbol x_0).
\]
The index $0$ denotes Kira's preferred source, $\boldsymbol x_q$ is the
candidate feature vector, and $h$ is the fitted tree. Lower targets
correspond to fewer recorded operations. We fix $\widehat y_0=0$ and
fit the remaining candidates with a weighted least-squares loss that
gives each training family equal weight. Validation selects tree
depth, minimum leaf population, target normalization and the required
predicted gain.

The selected tree has 21 nodes, depth four and at least two training
examples in each leaf. Its training set contains 68 candidate
continuations in 24 states from six families, while validation contains
30 candidates in nine states from three families. All 98 continuations
match their exact reference reductions. The tree changes nine training
choices and three validation choices, with family-weighted continuation
count ratios of 0.998602 and 0.999967 to the native choices. These
ratios measure individual choices followed by native continuation;
the next test must determine the effect of using the model repeatedly.

The model chooses the lowest predicted score and reports
\[
 P(q\mid S)=\frac{\exp(-\widehat y_q)}
                 {\sum_{q'}\exp(-\widehat y_{q'})}.
\]
This quantity expresses a relative model score, with no calibration
claim. Kira evaluates each score directly in compiled code. Groups with two rows or more than
32 rows keep the native source, which also avoids feature construction
for those groups. We preserve the original neural checkpoint and both
earlier linear fits as separate artifacts, and hold the final tree
fixed during reserved testing.

The seed implementation also reads its unchanged weight archive once
per process and retains at most two sets of family and candidate
embeddings within a workflow. A new target batch still receives its
own target features and scores. Independent checks compare every
probability with the original implementation and compare native tree
scores and choices with the Python model. These changes reduce
repeated loading and feature preparation while preserving the tested
seed predictions.

\subsection{Complete timing and reference checks}
Every final trial starts a new process and includes input preparation,
model loading and inference, native IBP generation and selection,
symbolic elimination, back substitution, exact coefficient comparison,
and all failed attempts and fallback reductions. Evidence compression
follows the stop time. These checks establish reduction to the prescribed reference basis;
they do not prove global minimality of the master set.
We report the initial reference calculation separately as common
benchmark preparation. Each policy receives the same prescribed
master basis and exact comparison result within a case.

All controls use the responsive Fermat shutdown from the preceding
section. We compare both the unchanged reduction executable with that
shutdown fix and an instrumented executable that keeps native row
choices. This separates the cost of measurement and integration from
the learned action. We also compare smaller deterministic regions and
matched deterministic seed subsets, since either could explain a gain
over conservative bounds.

Each trial uses one calculation thread, a bounded scratch directory
and limits on elapsed time and process-tree memory. Production trials
have 128 MiB of scratch space, a 2,400 MiB solver address-space limit, and a 2,100 MiB solver process-tree resident-memory limit. The enclosing worker has a 2,800 MiB address-space limit and a 2,600 MiB process-tree resident-memory limit. A single common-reference calculation
can use 256 MiB while retaining the same symbolic bounds, and the
storage ledger includes additional monitoring headroom. Final comparisons
use paired cases and report completion counts alongside time ratios.
Repeated runs reduce timing noise within a case, while different
kinematic instances remain distinct from different topologies.
The reserved results do not select the policy or its thresholds.
